% !TEX program = pdflatex
\documentclass[11pt]{article}

% ---------- Packages ----------
\usepackage[margin=1in]{geometry}
\usepackage{amsmath, amssymb, amsthm, mathtools}
\usepackage{siunitx}
\usepackage{physics}
\usepackage{booktabs}
\usepackage{graphicx}
\usepackage{xcolor}
\usepackage{microtype}
\usepackage{enumitem}
\usepackage{hyperref}
\usepackage[nameinlink,noabbrev]{cleveref}
\usepackage{algorithm}
\usepackage[noend]{algpseudocode}
\usepackage{listings}

% ---------- Listings setup ----------
\definecolor{lstbg}{RGB}{248,248,248}
\definecolor{lstkw}{RGB}{0,92,184}
\definecolor{lstid}{RGB}{100,0,140}
\definecolor{lstcm}{gray}{0.35}
\lstdefinestyle{code}{
  backgroundcolor=\color{lstbg},
  basicstyle=\ttfamily\small,
  keywordstyle=\color{lstkw}\bfseries,
  identifierstyle=\color{lstid},
  commentstyle=\color{lstcm}\itshape,
  showstringspaces=false,
  frame=single,
  framerule=0.2pt,
  rulecolor=\color{black!40},
  breaklines=true,
  upquote=true,
  columns=fullflexible
}

% ---------- Theorem-like ----------
\newtheorem{definition}{Definition}
\newtheorem{theorem}{Theorem}
\newtheorem{lemma}{Lemma}
\newtheorem{proposition}{Proposition}
\theoremstyle{remark}
\newtheorem{remark}{Remark}

% ---------- Macros ----------
\newcommand{\R}{\mathbb{R}}
\newcommand{\E}{\mathbb{E}}
\newcommand{\Var}{\mathrm{Var}}
\newcommand{\Cov}{\mathrm{Cov}}
\newcommand{\dd}{\,\mathrm{d}}
\newcommand{\Dt}{\,\Delta t}
\DeclareMathOperator*{\argmin}{arg\,min}
\DeclareMathOperator*{\argmax}{arg\,max}

% ---------- Title ----------
\title{A Falsification-First Protocol for ENG8-Style LENR\\with Deconvolved Calorimetry and Prime-Coherence Analysis}
\author{Ryan Van Gelder \& Tyler Van Osdol}
\date{\today}

\begin{document}
\maketitle

\begin{abstract}
We formalize a preregistered, falsification-first protocol to test claims of excess heat in ENG8-style low-energy nuclear reaction (LENR) devices. The method couples (i) rigorous energy accounting with deconvolved flow calorimetry, (ii) conservative chemical/physical energy bounds, (iii) model comparison with regularized priors penalizing extraordinary effects, and (iv) an optional prime-lattice coherence predictor that probes field-mediated mechanisms without entangling them with the primary energy test. We provide decision rules, uncertainty handling, blind-analysis procedures, and a reference Python analysis pipeline.
\end{abstract}

\tableofcontents

\section{Introduction \\ 
Scope and Philosophy}
Extraordinary heat-production claims require decision criteria that cannot be satisfied by calibration drift, unmetered RF power, or chemical storage. Our approach prioritizes falsifiability and replication. The protocol is architected in layers: (A) black-box energy test; (B) internal state modeling (optional); (C) prime-coherence mechanism probe (optional). Only (A) is required to declare excess heat.

\subsection*{Notation}
We denote the device-under-test (DUT) thermal state by temperature $T(t)$, the calorimeter impulse response by $h_{\mathrm{cal}}(t)$, instantaneous electrical input power by $P_{\mathrm{in}}^{\mathrm{el}}(t)$, deconvolved calorimetric output power by $P_{\mathrm{out}}^{\mathrm{cal}}(t)$, parasitic losses by $P_{\mathrm{loss}}(t)$, and reaction power by $P_{\mathrm{rxn}}(t)$.

\section{Energy Balance with Deconvolved Calorimetry}
\label{sec:energy}
In a closed energy accounting, the DUT satisfies
\begin{equation}
C_{\mathrm{th}}(T)\,\dot T(t) \;=\; P_{\mathrm{in}}^{\mathrm{el}}(t) \; - \; P_{\mathrm{out}}^{\mathrm{cal}}(t) \; - \; P_{\mathrm{loss}}(T,t) \; + \; P_{\mathrm{rxn}}(t) \; + \varepsilon(t).
\label{eq:balance}
\end{equation}
Rather than relying on raw $\Delta T(t)=T_{\mathrm{out}}-T_{\mathrm{in}}$, we model the calorimeter as linear time-invariant (LTI):
\begin{equation}
\Delta T(t) = (h_{\mathrm{cal}} * P_{\mathrm{true}})(t) + n(t), \quad P_{\mathrm{out}}^{\mathrm{cal}}(t) = \mathrm{deconv}\big[\Delta T(t),\,h_{\mathrm{cal}}(t);\,\lambda\big]\;\dot{m}(t)\,c_p.
\end{equation}
Here $\lambda$ is a regularization hyperparameter. We estimate $h_{\mathrm{cal}}$ via calibrated heater steps or multisine inputs, then perform frequency-domain Tikhonov (or Wiener) deconvolution; see \Cref{alg:ident,alg:deconv}.

\subsection{Excess Energy Criterion}
Define the integrated reaction energy
\begin{equation}
E_{\mathrm{rxn}}\;=\;\int_{t_0}^{t_1}\!\Big(P_{\mathrm{out}}^{\mathrm{cal}}(t) + P_{\mathrm{loss}}(t) - P_{\mathrm{in}}^{\mathrm{el}}(t)\Big)\,\dd t.
\end{equation}
Let $E_{\mathrm{chem}}^{\max}$ be a conservative bound on all chemical/physical storage (\Cref{sec:chem}). With total uncertainty $\sigma_E$, we require
\begin{equation}
E_{\mathrm{rxn}} - E_{\mathrm{chem}}^{\max} > 5\,\sigma_E.
\label{eq:gate}
\end{equation}
\paragraph{Uncertainty} We propagate sensor and model uncertainties using a bootstrap over $(v,i,T,\dot m)$ and over $\lambda$ (treated as a nuisance parameter). Transfer standard uncertainties for the electrical chain and the heater calibration are included.

\section{Calorimeter Identification and Deconvolution}
\label{sec:calid}
\begin{algorithm}[H]
\caption{Calorimeter impulse-response identification (frequency-domain)}
\label{alg:ident}
\begin{algorithmic}[1]
\Require Cal heater input $P_{\mathrm{cal}}[k]$, measured $\Delta T_{\mathrm{cal}}[k]$, sampling rate $f_s$, regularizer $\rho$
\State Compute $\mathcal{F}\{P\}(\omega)$ and $\mathcal{F}\{\Delta T\}(\omega)$ via FFT
\State $H(\omega) \gets \dfrac{\mathcal{F}\{\Delta T\}(\omega)}{\mathcal{F}\{P\}(\omega) + \rho\,\max_\nu |\mathcal{F}\{P\}(\nu)|}$ \Comment{Tikhonov-like stabilization}
\State $h_{\mathrm{cal}}[k] \gets \mathcal{F}^{-1}\{H(\omega)\}$; zero small negative-time wrap and truncate tail
\State \Return $h_{\mathrm{cal}}$
\end{algorithmic}
\end{algorithm}

\begin{algorithm}[H]
\caption{Deconvolution of heat flow}
\label{alg:deconv}
\begin{algorithmic}[1]
\Require $\Delta T[k]$, $h_{\mathrm{cal}}[k]$, $f_s$, regularizer $\lambda$
\State Compute $H(\omega)$ and $\mathcal{F}\{\Delta T\}(\omega)$
\State $P(\omega) \gets \dfrac{H^*(\omega)\,\mathcal{F}\{\Delta T\}(\omega)}{|H(\omega)|^2 + \lambda\,\max_\nu |H(\nu)|^2}$
\State $\widehat P_{\mathrm{out}}^{\mathrm{cal}}[k] \gets \mathcal{F}^{-1}\{P(\omega)\}$
\State \Return $\widehat P_{\mathrm{out}}^{\mathrm{cal}}[k]$
\end{algorithmic}
\end{algorithm}

\section{Chemical/Physical Energy Bound}
\label{sec:chem}
We include (i) bulk and double-layer electrostatic storage, (ii) phase transitions and thermal storage, (iii) gas recombination and adsorption/desorption.
\begin{align}
E_{\mathrm{chem}}^{\max} &= \tfrac{1}{2}C_{\mathrm{eff}}V_{\max}^2 + \tfrac{1}{2}C_{\mathrm{dl}}V_{\max}^2 \\
&\quad + E_{\mathrm{phase}}^{\max} + E_{\mathrm{gas\,recomb}}^{\max} + E_{\mathrm{ads}}^{\max}.
\end{align}
All inputs and uncertainties are pre-specified in the registry; see \Cref{tab:chem}.

\begin{table}[h]
\centering
\caption{Template for conservative chemical/physical energy bound (fill before data collection).}
\label{tab:chem}
\begin{tabular}{lccc}
\toprule
Component & Symbol & Value (CI) & Notes \\
\midrule
Bulk capacitance & $C_{\mathrm{eff}}$ & (\si{F}) & from EIS or datasheet \\
Double-layer cap. & $C_{\mathrm{dl}}$ & (\si{F}) & from EIS \\
Voltage span & $V_{\max}$ & (\si{V}) & preregistered operating range \\
Phase/latent heat & $E_{\mathrm{phase}}^{\max}$ & (\si{J}) & porosity/PCM/storage \\
Gas recombination & $E_{\mathrm{gas\,recomb}}^{\max}$ & (\si{J}) & H$_2$ + O$_2$ worst-case \\
Adsorption/desorption & $E_{\mathrm{ads}}^{\max}$ & (\si{J}) & surface area \& heats \\
\bottomrule
\end{tabular}
\end{table}

\section{Instrumentation and Loophole Closure}
We require four-quadrant wattmetry at DUT terminals, MHz-bandwidth logging of $v,i$, near-field $E/B$ probes, star-ground, RF chokes, and shield current monitors. Flow calorimetry uses redundant temperature sensors and mass-flow. Radiation channels (thermal neutron and $\gamma$) set upper bounds on particle power.

\section{Model Comparison and Priors}
\label{sec:model}
We compare $M_0$ (no reaction) versus $M_1$ (reaction power depends on loading and fields). Let $\theta(t)\in[0,1]$ denote a hydrogen loading proxy and $\mathcal{M}(t)=\sum_m w_m |a_m(t)|^2$ a modal field-energy proxy.
\begin{align}
M_0:&\quad y(t) = \beta_0 + \beta_1\,\theta(t) + \epsilon(t), \\
M_1:&\quad y(t) = \beta_0' + \beta_1'\,\theta(t) + \beta_2'\,\theta(t)^2\,\mathcal{M}(t) + \epsilon'(t),
\end{align}
where $y(t)=P_{\mathrm{out}}^{\mathrm{cal}}(t) - P_{\mathrm{in}}^{\mathrm{el}}(t)$ (optionally minus a static $P_{\mathrm{loss}}$ baseline).

\paragraph{Priors.} We regularize extraordinary effects: $\beta_2' \sim \mathrm{HalfNormal}(\sigma=0.1)$, while chemical baselines are weakly-informative (e.g., $\beta_0',\beta_1'\sim \mathcal{N}(0,1)$). Evidence is summarized by BIC or Bayes factors with a preregistered threshold $\Delta\mathrm{BIC}>10$.

\section{Prime-Coherence Mechanism Probe (Optional)}
Let $\{\varphi_p\}$ be prime-indexed basis filters derived from measured RF/near-field data. Define complex projections $S_p(t)=\langle s(t),\varphi_p\rangle$ over sliding windows, with amplitude $\hat\lambda_p=|S_p|$ and phase $\hat\phi_p=\arg S_p$. The coherence statistic is
\begin{equation}
\mathcal{C}(t) = \Big|\sum_{p} \hat\lambda_p(t)\,e^{\mathrm{i}\hat\phi_p(t)}\Big|^2.
\end{equation}
A significant partial correlation $\mathrm{corr}\big(P_{\mathrm{rxn}},\mathcal{C}\,|\,P_{\mathrm{in}}^{\mathrm{el}},T\big)>0$ across replicated runs supports a field-coherent mechanism; failure to observe it falsifies that mechanism without affecting the primary energy test.

\section{Blind Analysis and Decision Rules}
\begin{enumerate}[label=\textbf{D\arabic*.}]
  \item Pre-register code and parameters; lock a git hash before any DUT/control labels are revealed.
  \item Analyze blinded IDs with the frozen pipeline; compute $E_{\mathrm{rxn}}$, $\sigma_E$, and test \eqref{eq:gate}.
  \item If the DUT passes energy and model gates while controls do not, replicate on independent hardware. Only then unblind and declare.
\end{enumerate}

\section{Materials \\ Characterization and By-Products}
Pre- and post-experiment SEM/EDS of electrodes; gas chromatography of effluent; isotopic analysis (H/D, Li, Ni, etc.) if excess heat is declared. Radiation non-detections are converted to upper bounds using detector efficiency and geometry.

\section{Reproducibility: Data Schema}
Each run stores synchronized streams and metadata as in \Cref{tab:schema}.
\begin{table}[h]
\centering
\caption{Minimal data schema per run.}
\label{tab:schema}
\begin{tabular}{ll}
\toprule
Path & Description \\
\midrule
\texttt{/timebase} & sampling rate $f_s$ and $t_0$ (UTC) \\
\texttt{/v\_i} & columns $[v,i]$ at $f_s$ \\
\texttt{/thermal} & columns $[T_{\mathrm{in}},T_{\mathrm{out}}]$ \\
\texttt{/flow} & column $[\dot m]$ \\
\texttt{/calibration} & $[P_{\mathrm{cal\_in}},\,\Delta T_{\mathrm{cal}}]$ \\
\texttt{/fields} (opt.) & $[E_x,E_y,E_z,B_x,B_y,B_z,V_{\mathrm{RF}},I_{\mathrm{RF}}]$ \\
\texttt{/radiation} (opt.) & count rates or spectra \\
\texttt{/meta} & sensor models, wiring, geometry, uncertainties \\
\bottomrule
\end{tabular}
\end{table}

\section{Safety Considerations}
Operate within electrical, pressure, and radiation safety envelopes; implement interlocks and logging for any excursions. Treat unknown gas evolution as potentially flammable.

\appendix
\section{Reference Python Analysis Pipeline}
\label{app:pipeline}
\lstset{style=code,language=Python}
\begin{lstlisting}
import numpy as np
from numpy.fft import rfft, irfft

# ---------- Utilities ----------
def sync_and_trim(*arrays):
    L = min(map(len, arrays))
    return [a[:L] for a in arrays]

def real_input_power(v, i, fs):
    v, i = sync_and_trim(v, i)
    p = v * i
    E_in = np.trapz(p, dx=1/fs)
    return p, E_in

# ---------- Calorimeter ID ----------
def estimate_h_cal(P_cal, dT_cal, fs, reg=1e-4, max_len=None):
    P_cal, dT_cal = sync_and_trim(P_cal, dT_cal)
    N = int(2**np.ceil(np.log2(len(P_cal))))
    P = rfft(np.pad(P_cal, (0, N - len(P_cal))))
    T = rfft(np.pad(dT_cal, (0, N - len(dT_cal))))
    H = T / (P + reg * np.max(np.abs(P)))
    h = irfft(H)
    h = np.real(h)
    if max_len is not None:
        h = h[:max_len]
    h[h < 0] = 0.0  # enforce causality softly
    return h

# ---------- Deconvolution ----------
def deconvolve_power(dT, h, fs, reg=1e-4):
    L = len(dT) + len(h) - 1
    N = int(2**np.ceil(np.log2(L)))
    H = rfft(np.pad(h, (0, N - len(h))))
    T = rfft(np.pad(dT, (0, N - len(dT))))
    Hc = np.conj(H)
    denom = Hc*H + reg*np.max(Hc*H)
    P_est = irfft((Hc * T) / denom)
    return np.real(P_est[:len(dT)])

# ---------- Loss Model ----------
def loss_power(T_out, T_env, params):
    sigmaA = params.get('sigmaA', 0.0)
    k_cond = params.get('k_cond', 0.0)
    return sigmaA*(T_out**4 - T_env**4) + k_cond*(T_out - T_env)

# ---------- Chemical Bound ----------
def chemical_energy_bound(spec):
    E_cap = 0.5 * spec.get('C_eff_F', 0.0) * (spec.get('V_span', 0.0)**2)
    E_dl  = 0.5 * spec.get('C_dl_F', 0.0) * (spec.get('V_span', 0.0)**2)
    E_phase = spec.get('E_phase_J', 0.0)
    E_gas   = spec.get('E_gas_recomb_J', 0.0)
    E_ads   = spec.get('E_adsorption_J', 0.0)
    return E_cap + E_dl + E_phase + E_gas + E_ads

# ---------- Features ----------
def build_features(theta, T, modal_energy):
    X0 = np.column_stack([np.ones_like(T), theta])
    X1 = np.column_stack([np.ones_like(T), theta, theta**2 * modal_energy])
    return X0, X1

def bic(y, yhat, k, n):
    rss = np.sum((y - yhat)**2)
    sigma2 = rss / n
    return n*np.log(sigma2 + 1e-30) + k*np.log(n)

def fit_ridge(X, y, lam=10.0):
    XT = X.T
    A = XT @ X + lam * np.eye(X.shape[1])
    b = XT @ y
    w = np.linalg.solve(A, b)
    yhat = X @ w
    return w, yhat

# ---------- Prime-Coherence (windowed) ----------
def compute_coherence(sensor, prime_basis):
    keys = ['E_x','E_y','E_z','B_x','B_y','B_z','V_RF','I_RF']
    S = np.column_stack([sensor[k] for k in keys if k in sensor])
    s = np.sum(S, axis=1).astype(np.complex128)
    projections = []
    for phi in prime_basis:
        Sp = np.vdot(phi, s) / len(s)
        projections.append(Sp)
    lam = np.abs(np.array(projections))
    ang = np.angle(np.array(projections))
    C = np.abs(np.sum(lam * np.exp(1j*ang)))**2
    return C, lam, ang

# ---------- Main ----------
def analyze_run(data, params):
    fs = data['fs']
    v, i = data['v'], data['i']
    T_in, T_out = data['T_in'], data['T_out']
    m_dot = data['m_dot']
    cp = params.get('cp_J_per_kgK', 4180.0)
    T_env = params.get('T_env_K', 298.15)

    p_in, E_in = real_input_power(v, i, fs)

    dT = T_out - T_in
    h = estimate_h_cal(data['P_cal_in'], data['dT_cal'], fs=data['fs_cal'],
                       reg=params.get('reg_id',1e-4), max_len=params.get('h_len'))
    P_out = deconvolve_power(dT, h, fs, reg=params.get('reg_deconv',1e-4)) * (m_dot * cp)

    P_loss = loss_power(T_out, T_env, params.get('loss', {}))

    P_rxn = P_out + P_loss - p_in
    E_rxn = np.trapz(P_rxn, dx=1/fs)

    sigma_P = np.std(P_rxn)
    sigma_E = sigma_P * (len(P_rxn)/fs)**0.5

    E_chem_max = chemical_energy_bound(params.get('chem_spec', {}))

    results = {
        'E_in_J': float(E_in),
        'E_rxn_J': float(E_rxn),
        'E_chem_max_J': float(E_chem_max),
        'sigma_E_J': float(sigma_E),
        'passes_energy_gate': bool((E_rxn - E_chem_max) > 5*sigma_E)
    }

    if all(k in data for k in ['theta','modal_energy']):
        y = P_out - p_in
        X0, X1 = build_features(data['theta'], (T_in+T_out)/2, data['modal_energy'])
        w0, y0 = fit_ridge(X0, y, lam=params.get('lam0', 10.0))
        w1, y1 = fit_ridge(X1, y, lam=params.get('lam1', 10.0))
        n = len(y)
        results.update({
            'BIC_M0': float(bic(y, y0, k=X0.shape[1], n=n)),
            'BIC_M1': float(bic(y, y1, k=X1.shape[1], n=n)),
            'delta_BIC': float(bic(y, y0, k=X0.shape[1], n=n) - bic(y, y1, k=X1.shape[1], n=n)),
            'supports_FIELD_term': bool((bic(y, y0, k=X0.shape[1], n=n) - bic(y, y1, k=X1.shape[1], n=n)) > 10)
        })

    if 'fields' in data and 'prime_basis' in data:
        C, lam, ang = compute_coherence(data['fields'], data['prime_basis'])
        results.update({'prime_C': float(C)})

    return results
\end{lstlisting}

\section{Calibration and Synthetic Dry-Run}
\subsection{Deconvolution Sanity Test}
\lstset{style=code,language=Python}
\begin{lstlisting}
# Generate a two-pole calorimeter and verify recovery
fs = 10.0
T = 3600  # seconds
N = int(fs*T)

# True impulse response: h(t) = a1*exp(-t/t1) + a2*exp(-t/t2)
import numpy as np

t = np.arange(N)/fs
h_true = 0.8*np.exp(-t/50.0) + 0.2*np.exp(-t/300.0)

# Input heat: PRBS steps
rng = np.random.default_rng(1)
steps = rng.choice([0, 10.0], size=N)
P_true = steps

# Observed dT = h * P + noise
from numpy.fft import rfft, irfft
L = len(P_true) + len(h_true) - 1
M = int(2**np.ceil(np.log2(L)))
H = rfft(np.pad(h_true, (0, M - len(h_true))))
P = rfft(np.pad(P_true, (0, M - len(P_true))))
Tspec = H*P

dT_noiseless = irfft(Tspec)[:len(P_true)]
dT = dT_noiseless + 0.01*rng.standard_normal(len(dT_noiseless))

# Estimate h and deconvolve
h_est = estimate_h_cal(P_true, dT, fs, reg=1e-4, max_len=int(8*fs*300))
P_rec = deconvolve_power(dT, h_est, fs, reg=1e-4)

# Check recovery error
err = np.linalg.norm(P_true - P_rec)/np.linalg.norm(P_true)
print('Relative L2 error:', err)
\end{lstlisting}

\section{Preregistration Form (Template)}
List git commit hash of analysis code; calorimeter regularization grid; sensor models and uncertainties; DUT/control labels; stop/declare thresholds; replication plan.

\section*{Acknowledgments}
We thank collaborators for feedback on identifiability, deconvolution, and blind-analysis design.
\nocite{*}
\bibliographystyle{unsrt}
\bibliography{references}
\end{document}
